% Бүлэг 1

\chapter{Онолын хэсэг}
\label{Chapter1}
\pagecolor{white}
%-------------------------------------------------------------------------------

\newcommand{\keyword}[1]{\textbf{#1}}
\newcommand{\tabhead}[1]{\textbf{#1}}
\newcommand{\code}[1]{\texttt{#1}}
\newcommand{\file}[1]{\texttt{\bfseries#1}}
\newcommand{\option}[1]{\texttt{\itshape#1}}

\definecolor{accentblue}{RGB}{52,152,219}
\definecolor{accentgreen}{RGB}{39,174,96}
\definecolor{accentpurple}{RGB}{142,68,173}
\definecolor{accentorange}{RGB}{230,126,34}
\definecolor{accentred}{RGB}{231,76,60}
\definecolor{accentteal}{RGB}{26,188,156}
\definecolor{accentyellow}{RGB}{241,196,15}
\definecolor{darknavy}{RGB}{44,62,80}
\definecolor{softgrey}{RGB}{236,240,241}
\definecolor{midgrey}{RGB}{189,195,199}
\definecolor{lightblue}{RGB}{174,214,241}

% Хүснэгтийн нэгдсэн стиль (бүх хүснэгт дээр ашиглана)
% \footnotesize нь body text-ээс жижиг хэмжээтэй
\newcommand{\tabletop}[1]{\rowcolor{darknavy}\textcolor{white}{\textbf{#1}}}

%-------------------------------------------------------------------------------

\section{Сэдэв сонгох үндэслэл}
Орчин үеийн нийгэмд хувийн санхүүгийн мэдлэг, ур чадварыг эзэмших нь иргэн бүрийн хувьд зайлшгүй шаардлага болсон. Гар утас, мобайл банкны үйлчилгээ хөгжсөнөөр иргэдийн өдөр тутмын гүйлгээний \textbf{тоо, давтамж эрс нэмэгдсэн} бөгөөд үүнтэй уялдан үүсэх санхүүгийн мэдээллийг гараар бүртгэх, ангилах, шинжлэх боломж улам бүр хязгаарлагдмал болж байна. Тиймээс \textbf{санхүүгийн мэдээллийг автоматаар нэгтгэн боловсруулах}, хэрэглэгчид тодорхой, ойлгомжтой хэлбэрээр хүргэх систем зайлшгүй шаардлагатай.\\

Энэхүү сэдвийг сонгох үндэслэлийг дараах хэдэн чухал хүчин зүйлээр тайлбарлаж болно:

\begin{itemize}
    \item [\textbf{1.}] \textbf{Иргэдийн санхүүгийн боловсролын түвшин харьцангуй сул.} S\&P Global FinLit судалгаагаар Монгол Улсын насанд хүрсэн хүн амын ердөө 30 орчим хувь нь санхүүгийн суурь ойлголтуудыг сайн ойлгодог \cite{mn_finlit_survey}. Энэ нь үндэсний эдийн засагт сөргөөр нөлөөлдөг бөгөөд хэт зарцуулалт, хуримтлал бага байх, өрийн дарамт өсөх зэрэг бодит асуудлуудыг бий болгож байна.

    \item [\textbf{2.}] \textbf{Дотоодын зах зээлд тохирсон шийдэл байхгүй.} Олон улсын Mint, YNAB, Cleo зэрэг шийдлүүд нь Монголын банкны экосистем, Монгол хэл, төгрөгийн валюттай ажиллах боломжгүй. Дотоодын банкны гар утасны программууд нь зөвхөн өөрийн дансны мэдээллийг харуулах хүрээтэй бөгөөд хувийн санхүүгийн \textbf{нэгдсэн, ухаалаг хяналт} хийх боломж хязгаарлагдмал.

    \item [\textbf{3.}] \textbf{Хиймэл оюун ухааны технологи бэлэн болсон.} Сүүлийн жилүүдэд том хэлний загвар (LLM)-ууд хүртээмжтэй болсон бөгөөд Google Gemini, OpenRouter зэрэг үнэгүй ба нээлттэй платформуудын тусламжтайгаар хувийн санхүүгийн \textbf{өндөр чанартай зөвлөмжийн систем} бий болгох техникийн бүх боломж нээгдсэн.

    \item [\textbf{4.}] \textbf{Олон банкны данс ашиглах хандлага түгээмэл болсон.} Нэг хэрэглэгч нь 2--3 банкны данс зэрэг ашиглах нь өнөөдөр түгээмэл практик болсон. Иймд эдгээр өгөгдлийг нэгтгэж нэг дороос харуулдаг \textbf{aggregator (нэгтгэгч) платформ} зах зээлийн бодит эрэлттэй болсон.

    \item [\textbf{5.}] \textbf{Гар утсанд суурилсан хэрэглээний боломж.} iOS, Android платформын хүртээмж 90 хувиас давсан өнөө үед хувийн санхүүг \textbf{бодит цаг хугацаанд гар утсаар хянах} боломж хамгийн тохиромжтой шийдэл болж байна.

    \item [\textbf{6.}] \textbf{Хувь хэрэглэгчийн талаас өөрийн туршлага.} Зохиогчийн хувьд олон банкны хуулгыг гараар нэгтгэх, Excel дээр оруулах ажил их хугацаа шаардаж, удаан үргэлжилсэн. Үүнийг автоматжуулах шаардлага нь сэдвийг сонгох \textbf{практик хүчин зүйл} болсон.
\end{itemize}

\textbf{Сэдвийн шинэлэг чанар.} Монголын зах зээл дээр хувийн санхүүг олон банкаас нэгтгэн, AI-аар автомат шинжилж, хэрэглэгчид хувийн зөвлөмж өгдөг \textbf{цогц шийдэл} одоогоор алга байна. Энэ нь дипломын ажлын онцлог, шинэлэг тал юм.

\section{Тулгамдсан асуудал ба шийдвэрлэх арга зам}
Хувийн санхүүгийн менежментийн өнөөгийн нөхцөлд хэрэглэгчдийн тулгардаг үндсэн асуудлууд болон энэхүү дипломын ажлаар санал болгож буй шийдлийг харьцуулан үзэв.

\begin{table}[ht]
\centering
\footnotesize
\renewcommand{\arraystretch}{1.35}
\rowcolors{2}{softgrey}{white}
\begin{tabular}{|p{4.4cm}|p{8cm}|}
\hline
\rowcolor{darknavy}
\textcolor{white}{\textbf{Тулгамдсан асуудал}} & \textcolor{white}{\textbf{FinTrack-ийн санал болгож буй шийдэл}} \\ \hline
Олон банкны мэдээлэл сарнисан байх & Бүх банкны хуулгыг нэг программд PDF / Excel / CSV форматаар оруулж, AI-аар автоматаар задлан, нэг дансны төвд харуулах. \\ \hline
Гар бүртгэл цаг их шаарддаг & Гүйлгээ нэмэх дэлгэц нь хэдхэн даралтаар бөглөгдөх, дансны үлдэгдэл шууд автоматаар шинэчлэгдэх дизайнтай. \\ \hline
Зарлагыг ангилахад хүндрэлтэй & Хоёр түвшний автомат ангилал: дүрэмд суурилсан танилт (rule-based) ба LLM-д суурилсан утга зүйн (semantic) ухаалаг танилт. \\ \hline
Зарцуулалтын хэв маяг ойлгомжгүй & Сар, улирал, жилээр статистик, дугуй / багана / гөлгөр муруй график, ангиллын задаргаа автоматаар гаргах. \\ \hline
Эрсдэлийн мэдээлэл хоцрогддог & Push мэдэгдлээр (Expo Push) төсөв хэтрэх, том дүнгийн гүйлгээ үүсэх үед бодит цагт сэрэмжлүүлэх. \\ \hline
Хувийн санхүүгийн зөвлөгөө байхгүй & 30 хоногийн контекст ашигладаг AI чатбот --- хэрэглэгчийн бодит өгөгдлийн үндсэн дээр зөвлөгөө өгөх. \\ \hline
Давтагдах төлбөр мартагдах & Сар тутмын тогтмол төлбөрийг (түрээс, интернэт, утас, даатгал) урьдчилан бүртгэж, эрт сануулга өгөх модуль. \\ \hline
Олон валютын хөрвөгдөл & MNT, USD, EUR, KRW, CNY, JPY, RUB, GBP гэсэн 8 валютыг автомат ханшаар хөрвүүлэх. \\ \hline
Аюулгүй байдлын эрсдэл & JWT баталгаажуулалт, bcrypt нууц үг хэшлэлт, Face ID / Touch ID биометрик нэвтрэлт, аудит лог IP хаягтай. \\ \hline
LLM-ийн газарзүйн хязгаарлалт & Олон нийлүүлэгчийн (Gemini $\to$ OpenRouter) автомат нөөц шилжилтийн архитектур ашиглан тасралтгүй ажиллагааг хангах. \\ \hline
\end{tabular}
\caption{Тулгамдсан асуудал ба санал болгож буй шийдэл}
\label{tab:problem-solution}
\end{table}

\textbf{Шийдвэрлэх ерөнхий арга зам.} Дээрх асуудлуудыг шийдвэрлэхийн тулд дараах техникийн стратегийг сонгосон:

\begin{itemize}
    \item [\textbf{1.}] \textbf{Дөрвөн давхаргат архитектур} --- танилцуулах (React Native), бизнес логик (Go + Gin), өгөгдөл (PostgreSQL + GORM), AI / гадаад үйлчилгээ гэсэн давхаргад тус бүр өөрийн хариуцлагын хүрээнд ажиллана.

    \item [\textbf{2.}] \textbf{Хуулга задлагч микросервис} --- Python + FastAPI + pdfplumber дээр суурилсан тусгай үйлчилгээ нь Монгол Улсын томоохон таван банкны (Хаан, Голомт, ХХБ, Хас, Төрийн) хуулгыг шууд унших чадвартай.

    \item [\textbf{3.}] \textbf{Олон нийлүүлэгчтэй AI давхарга} --- Google Gemini 2.5 Flash-ийг үндсэн, OpenRouter (Llama 3.3, Qwen 2.5, Mistral, DeepSeek)-ийг нөөц болгон ашигладаг төв давхарга.

    \item [\textbf{4.}] \textbf{Контекстэд тулгуурласан AI зөвлөгөө} --- хэрэглэгчийн сүүлийн 30 хоногийн гүйлгээг JSON хэлбэрээр AI-ийн системийн зааварчилгаанд оруулж, хувийн нөхцөлд тохирсон зөвлөмж олж авна.

    \item [\textbf{5.}] \textbf{Component-based UI} --- React Native + NativeWind v4 ашиглан дахин ашиглаж болох жижиг UI components-оор нийт 20+ дэлгэцийг бүтээх.
\end{itemize}

Эдгээр шийдлийг хослуулснаар хэрэглэгч \textbf{нэг апп дотроосоо} санхүүгийн өгөгдлөө нэгтгэх, шинжлэх, AI-аас зөвлөгөө авах, эрсдэлийн мэдэгдэл хүлээж авах боломжийг олох болно.

\section{Судалгааны ажлын зорилго, зорилт}
\textbf{Судалгааны ажлын зорилго} нь хувь хүний санхүүгийн мэдээллийг нэгтгэн удирдах, гүйлгээний өгөгдөлд шинжилгээ хийх, улмаар хэрэглэгчийн шийдвэр гаргалтад дэмжлэг үзүүлэх ухаалаг зөвлөмжийн системийн загвар, хэрэгжилтийг боловсруулахад оршино.\\

Дээрх зорилгыг хэрэгжүүлэхийн тулд дараах \textbf{зорилтууд}-ыг дэвшүүлэв:
\begin{itemize}
    \item [1.] Хувийн санхүүгийн менежментийн өнөөгийн байдал, хэрэглэгчийн тулгамдсан асуудлыг судлах;
    \item [2.] Системд тавигдах функциональ болон функциональ бус шаардлагыг тодорхойлох;
    \item [3.] Гүйлгээ бүртгэл, ангилал, шинжилгээ, зөвлөмжийн үндсэн урсгал бүхий системийн зохиомж боловсруулах;
    \item [4.] Банкны хуулгад тулгуурласан автомат боловсруулалт, шинжилгээний механизмыг хэрэгжүүлэх;
    \item [5.] Хэрэглэгчийн нөхцөлд нийцсэн AI зөвлөмж болон чат харилцааны боломжийг нэвтрүүлэх;
    \item [6.] Боловсруулсан шийдлийн үр дүн, хэрэглээний ач холбогдлыг үнэлж дүгнэх.
\end{itemize}

\section{Хувийн санхүү ба түүнийг автоматжуулах хэрэгцээ}
\textbf{Хувийн санхүү} гэдэг нь хувь хүн өөрийн орлого, зарлага, хадгаламж, хөрөнгө оруулалт болон өр төлбөрөө урт хугацааны санхүүгийн зорилготойгоо уялдуулан төлөвлөж, үр ашигтайгаар удирдах үйл явц юм. Санхүүгээ зөв зохион байгуулснаар хүн өдөр тутмын хэрэглээгээ оновчтой хянах, ирээдүйн зорилгодоо хуримтлал үүсгэх, цаашлаад санхүүгийн тогтвортой байдалд хүрэх боломжтой болдог.\\

Хувийн санхүүгийн зөв менежмент нь богино хугацаанд өдөр тутмын зардлаа хэмнэх, дунд хугацаанд тодорхой зорилгод зориулсан хадгаламж бүрдүүлэх, урт хугацаанд тэтгэвэр, хөрөнгө оруулалт зэрэг ирээдүйн хэрэгцээг хангах үндэс суурийг бүрдүүлдэг.\\

Орчин үеийн хувь хүний санхүүгийн тогтолцоо нь дараах үндсэн чиглэлүүдээс бүрдэнэ. Үүнд:
\begin{itemize}
    \item [•] \textbf{Орлогын удирдлага.} Цалин хөлс, фриланс орлого, бизнесийн ашиг, хөрөнгийн өгөөж зэрэг бүх төрлийн орлогыг бүртгэж, төлөвлөх.
    \item [•] \textbf{Зарлагын төлөвлөлт.} Хүнс, тээвэр, хэрэглээний төлбөр, эрүүл мэнд, боловсрол, зугаа цэнгэл зэрэг ангиллаар зарлагаа хянах.
    \item [•] \textbf{Хадгаламж.} Яаралтай хэрэгцээний сан бүрдүүлэх, мөн орон сууц, автомашин, боловсрол зэрэг зорилтот хуримтлал үүсгэх.
    \item [•] \textbf{Хөрөнгө оруулалт.} Хувьцаа, бонд, хадгаламжийн бүтээгдэхүүн, үл хөдлөх хөрөнгө зэрэгт хөрөнгө байршуулах замаар хөрөнгийн өсөлтийг хангах.
    \item [•] \textbf{Өр, зээлийн удирдлага.} Хэрэглээний болон ипотекийн зээл, кредит картын хэрэглээг зохистой түвшинд удирдах.
    \item [•] \textbf{Даатгалын хамгаалалт.} Эрүүл мэнд, амь нас, эд хөрөнгийн эрсдэлийг бууруулах зорилгоор даатгал ашиглах.
    \item [•] \textbf{Татварын төлөвлөлт.} Татварын хөнгөлөлт, чөлөөлөлт болон холбогдох боломжуудыг зөв ашиглах.
\end{itemize}

Сүүлийн жилүүдэд дижитал технологийн хөгжилтэй холбоотойгоор хувийн санхүүг автоматжуулах хэрэгцээ эрчимтэй нэмэгдэж байна. Гар аргаар зарлага бүртгэх, төсөв тооцоолох нь цаг хугацаа их шаарддаг төдийгүй алдаа гарах магадлал өндөр байдаг. Харин автоматжуулсан систем ашигласнаар хэрэглэгчийн орлого, зарлагыг бодит цаг хугацаанд ангилж, дүн шинжилгээ хийхээс гадна санхүүгийн төлөвлөгөө боловсруулах, хэт зарцуулалтаас сэргийлэх, хуримтлал үүсгэхэд дэмжлэг үзүүлэх боломж бүрддэг.

\subsection{Хүн амын санхүүгийн боловсролын өнөөгийн байдал}
2024 онд хийгдсэн S\&P Global FinLit судалгаагаар дэлхийн насанд хүрсэн хүн амын зөвхөн \textbf{33\%} нь хүү, инфляц, эрсдэлийн хуваарилалт зэрэг санхүүгийн суурь ойлголтуудыг хангалттай түвшинд ойлгодог болохыг тогтоосон \cite{split_finlit, worldbank_finlit}. Монгол Улсын хувьд энэхүү үзүүлэлт ойролцоогоор \textbf{30\%} байгаа бөгөөд ялангуяа залуу үеийн дунд хувийн санхүүгээ зөв удирдах мэдлэг, ур чадвар харьцангуй сул түвшинд байна \cite{mn_finlit_survey}.\\

Нийтийн санхүүгийн боловсролыг нэмэгдүүлэх, иргэдийн санхүүгийн мэдлэгийг сайжруулах асуудлыг судалсан олон улсын болон дотоодын судалгааны ажлууд \cite{mn_finlit, mn_personalfin} энэ чиглэлийн ач холбогдлыг онцлон тэмдэглэсэн байдаг. Судалгааны үр дүнгээс харахад иргэд өдөр тутмын санхүүгээ удирдах явцад дараах нийтлэг хүндрэлүүдтэй тулгарч байна. Үүнд:
\begin{itemize}
    \item [•] Сар бүрийн орлого, зарлагын үлдэгдэл тодорхой бус байх,
    \item [•] Ямар төрлийн зардал хамгийн өндөр хувийг эзэлж байгааг тодорхойлох боломж хязгаарлагдмал байх,
    \item [•] Давтамжтай төлбөрүүд (захиалгат үйлчилгээ, зээлийн эргэн төлөлт гэх мэт) хяналтгүй өсөх,
    \item [•] Хадгаламж болон санхүүгийн зорилгод хүрэх тодорхой төлөвлөгөө дутмаг байх,
    \item [•] Орлого, зарлагын харьцаа болон хуримтлалын түвшин (savings rate)-ийг тооцоолох мэдлэг, дадал хангалтгүй байх,
    \item [•] Олон банкны данс ашигладаг тохиолдолд нийт хөрөнгийн үлдэгдлийг нэг дороос бодит цаг хугацаанд хянах боломжгүй байх зэрэг асуудлууд түгээмэл ажиглагдаж байна.
\end{itemize}

Эдгээр асуудлууд нь хувь хүний санхүүгийн сахилга бат, хуримтлалын түвшин, цаашлаад урт хугацааны санхүүгийн тогтвортой байдалд сөргөөр нөлөөлөх эрсдэлтэй юм.

\subsection{Уламжлалт банкны системийн хязгаарлалт}
Монгол Улсад үйл ажиллагаа явуулж буй арилжааны банкууд, тухайлбал Худалдаа Хөгжлийн Банк, Голомт банк, Хаан банк зэрэг байгууллагууд өөрсдийн мобайл болон веб үйлчилгээг хэрэглэгчдэд санал болгодог \cite{mn_bankstat, mn_monbank}. Эдгээр системүүд нь дансны үлдэгдэл харах, гүйлгээ хийх, төлбөр төлөх зэрэг үндсэн банкны үйлчилгээг хангадаг боловч хэрэглэгчийн хувийн санхүүгийн иж бүрэн удирдлага, шинжилгээний хэрэгцээг бүрэн хангаж чаддаггүй.\\

Өнөө үед хэрэглэгчид олон банкны данс зэрэг ашиглах нь түгээмэл болсон ч тухайн банкны систем нь зөвхөн өөрийн байгууллагын мэдээллийг харуулдаг онцлогтой байдаг. Үүнтэй холбоотойгоор дараах хязгаарлалтууд үүсдэг. Үүнд:
\begin{itemize}
    \item [•] Хэрэглэгчийн бүх банкны дансны нийт хөрөнгө, үлдэгдлийг нэг дороос харах боломж хязгаарлагдмал,
    \item [•] Гүйлгээний автомат ангилал (categorization) дутмаг бөгөөд зарлагыг "хүнс", "тээвэр", "зугаа цэнгэл" зэрэг төрлөөр автоматаар ялгах боломж хангалтгүй,
    \item [•] Хэрэглэгчийн зарцуулалтын хэв маяг, санхүүгийн зан төлөвийн өөрчлөлт болон чиг хандлагыг (trend analysis) тодорхойлох функц ховор,
    \item [•] Хиймэл оюун ухаанд суурилсан санхүүгийн зөвлөмж, эрсдэлийн анхааруулга, төсвийн санал зэрэг дэвшилтэт боломжууд бараг хөгжөөгүй,
    \item [•] Орлого, тогтмол зарлага болон хуримтлалын мэдээлэлд үндэслэн хэрэглэгчийн чөлөөт мөнгөн урсгал (cash flow)-ыг тооцоолж, дүн шинжилгээ хийх боломж хязгаарлагдмал байдаг.
\end{itemize}

Иймээс уламжлалт банкны системүүд нь үндсэн гүйлгээний үйлчилгээг хангаж байгаа боловч хэрэглэгчийн санхүүгийн зан төлөвийг иж бүрэн шинжлэх, ухаалаг зөвлөмж өгөх, олон эх үүсвэрийн мэдээллийг нэгтгэн боловсруулах зэрэг орчин үеийн хэрэгцээг бүрэн хангаж чадахгүй байна.

\begin{table}[ht]
\centering
\footnotesize
\renewcommand{\arraystretch}{1.25}
\rowcolors{2}{softgrey}{white}
\begin{tabular}{|p{5cm}|c|c|}
\hline
\rowcolor{darknavy}
\textcolor{white}{\textbf{Боломж, шинж чанар}} & \textcolor{white}{\textbf{Банкны программ}} & \textcolor{white}{\textbf{AI санхүүгийн систем}} \\ \hline
Олон банкны данс нэгтгэх       & $\times$    & $\checkmark$ \\ \hline
Автомат ангилал                & $\times$    & $\checkmark$ \\ \hline
AI зөвлөмж                     & $\times$    & $\checkmark$ \\ \hline
Сарын төсөв удирдах            & Зарим       & $\checkmark$ \\ \hline
Давтагдах төлбөрийн анхаарал   & $\times$    & $\checkmark$ \\ \hline
Чат хэлбэрээр зөвлөгөө         & $\times$    & $\checkmark$ \\ \hline
Push мэдэгдэл (дүгнэлт)        & $\times$    & $\checkmark$ \\ \hline
Хэрэглэгчийн өгөгдөл цуглуулагч & Зөвхөн өөрийнх & Хэрэглэгч өөрөө сонгоно \\ \hline
Гүйлгээний чиг хандлагын шинжилгээ & $\times$ & $\checkmark$ \\ \hline
Хадгаламжийн зорилгот тооцоо   & $\times$    & $\checkmark$ \\ \hline
Хэрэглэгчийн зан төлөв сурах   & $\times$    & $\checkmark$ \\ \hline
\end{tabular}
\caption{Уламжлалт банкны программ ба AI-д суурилсан санхүүгийн системийн харьцуулалт}
\label{tab:fin-app-cmp}
\end{table}

\subsection{Дэлхий дахинд хэрэглэгдэж буй ижил төрлийн шийдлүүд}
Сүүлийн жилүүдэд олон улсад хувийн санхүүгийн \emph{нэгтгэгч} төрлийн программ хангамжууд эрчимтэй хөгжиж, өргөн хэрэглээнд нэвтэрч байна. Эдгээр системүүд нь Open Banking API ашиглан хэрэглэгчийн банкны мэдээллийг автоматаар нэгтгэж, орлого, зарлага, хадгаламж болон хөрөнгийн мэдээллийг нэг дороос удирдах боломжийг бүрдүүлдэг. Мөн хэрэглэгчдэд санхүүгийн нэгдсэн хяналтын самбар, зарлагын ангилал, төсөв төлөвлөлт, санхүүгийн шинжилгээ зэрэг ухаалаг үйлчилгээг санал болгодог.\\

Гадаадын ижил төрлийн системүүд нь хиймэл оюун ухаан болон өгөгдлийн шинжилгээнд тулгуурлан хэрэглэгчийн зарцуулалтын хэв маягийг тодорхойлох, ирээдүйн мөнгөн урсгалыг таамаглах, хэт зарцуулалтаас урьдчилан сэргийлэх зөвлөмж өгөх зэрэг дэвшилтэт боломжуудыг агуулдаг.

\begin{table}[ht]
\centering
\footnotesize
\renewcommand{\arraystretch}{1.25}
\rowcolors{2}{softgrey}{white}
\begin{tabular}{|p{2.4cm}|p{2.6cm}|p{3cm}|p{4cm}|}
\hline
\rowcolor{darknavy}
\textcolor{white}{\textbf{Программ}} & \textcolor{white}{\textbf{Гарал}} & \textcolor{white}{\textbf{Гол давуу тал}} & \textcolor{white}{\textbf{Сул тал, дутагдал}} \\ \hline
Mint            & АНУ          & Олон банк нэгтгэх, төсөв        & 2024 онд хаагдсан, AI зөвлөмж сул \\ \hline
YNAB            & АНУ          & Бариу төсвийн арга              & Зөвхөн төлбөртэй, AI байхгүй \\ \hline
Monarch Money   & АНУ          & Шинэ үеийн UI ба UX                & Open Banking-гүй улсад хязгаарлагдмал \\ \hline
Wallet (BudgetBakers) & ЕХ      & Олон валют, олон улс            & AI зөвлөмж сул \\ \hline
Spendee         & Чех          & Гэр бүлийн төсөв                & Нэг хэрэглэгч төвлөрсөн \\ \hline
Toshl           & Словени      & Дизайн, тоглоомжуулалт          & Үнэ өндөр \\ \hline
\textbf{FinTrack (миний ажил)} & \textbf{Монгол} & \textbf{LLM AI чат + банкны хуулга шинжилгээ} & \textbf{Open Banking интеграц нэмэх боломжтой} \\ \hline
\end{tabular}
\caption{Дэлхийн зах зээл дэх алдартай шийдлүүд}
\label{tab:competitors}
\end{table}

\subsection{Функциональ боломжуудын дэлгэрэнгүй харьцуулалт}
Дээрх ерөнхий тоймыг гүнзгийрүүлэх үүднээс, шилдэг ижил төстэй шийдлүүдийг функциональ боломжоор нь нэг бүрчлэн харьцуулсан хүснэгтийг доор үзүүлэв. Тэмдэглэгээ: $\checkmark$ --- бүрэн дэмждэг, $\sim$ --- хэсэгчлэн / нэмэлт төлбөртэй, $\times$ --- дэмждэггүй.

\begin{table}[ht]
\centering
\scriptsize
\renewcommand{\arraystretch}{1.2}
\rowcolors{2}{softgrey}{white}
\begin{tabular}{|p{4.0cm}|c|c|c|c|c|c|}
\hline
\rowcolor{darknavy}
\textcolor{white}{\textbf{Боломж / Функц}} &
\textcolor{white}{\textbf{Mint}} &
\textcolor{white}{\textbf{YNAB}} &
\textcolor{white}{\textbf{Monarch}} &
\textcolor{white}{\textbf{Wallet}} &
\textcolor{white}{\textbf{Money Lover}} &
\textcolor{white}{\textbf{\textbf{FinTrack}}} \\ \hline
Олон банкны данс нэгтгэх        & $\checkmark$ & $\checkmark$ & $\checkmark$ & $\checkmark$ & $\sim$        & $\checkmark$ \\ \hline
Гар утсанд шууд ажиллах         & $\checkmark$ & $\checkmark$ & $\checkmark$ & $\checkmark$ & $\checkmark$  & $\checkmark$ \\ \hline
Сарын төсөв                     & $\checkmark$ & $\checkmark$ & $\checkmark$ & $\checkmark$ & $\checkmark$  & $\checkmark$ \\ \hline
Автомат ангилал                 & $\checkmark$ & $\sim$       & $\checkmark$ & $\sim$       & $\sim$        & $\checkmark$ \\ \hline
Олон валют                      & $\times$     & $\sim$       & $\sim$       & $\checkmark$ & $\checkmark$  & $\checkmark$ \\ \hline
Давтагдах төлбөрийн дохио       & $\sim$       & $\checkmark$ & $\checkmark$ & $\sim$       & $\checkmark$  & $\checkmark$ \\ \hline
Push мэдэгдэл                   & $\checkmark$ & $\checkmark$ & $\checkmark$ & $\checkmark$ & $\checkmark$  & $\checkmark$ \\ \hline
\textbf{LLM AI чат зөвлөгөө}    & $\times$     & $\times$     & $\sim$       & $\times$     & $\times$      & $\checkmark$ \\ \hline
\textbf{Банкны хуулгыг AI задлах} & $\times$   & $\times$     & $\times$     & $\sim$       & $\times$      & $\checkmark$ \\ \hline
\textbf{Монгол хэлний дэмжлэг}  & $\times$     & $\times$     & $\times$     & $\sim$       & $\sim$        & $\checkmark$ \\ \hline
\textbf{Монгол банкны (ХХБ/Голомт/Хаан) форматын дэмжлэг} & $\times$ & $\times$ & $\times$ & $\times$ & $\times$ & $\checkmark$ \\ \hline
Үнэгүй ашиглах                  & -- (хаагдсан) & $\times$     & $\times$     & $\sim$       & $\sim$        & $\checkmark$ \\ \hline
Open Banking интеграц           & $\checkmark$ & $\checkmark$ & $\checkmark$ & $\checkmark$ & $\sim$        & $\sim$(төлөвлөгөөтэй) \\ \hline
Хөрөнгө оруулалт хянах          & $\checkmark$ & $\times$     & $\checkmark$ & $\times$     & $\times$      & $\sim$(төлөвлөгөөтэй) \\ \hline
Өгөгдлийг өөрөө байршуулах (self-host) & $\times$ & $\times$ & $\times$  & $\times$     & $\times$      & $\checkmark$ \\ \hline
\end{tabular}
\caption{Ижил төстэй шийдлүүдийг функциональ боломжоор харьцуулсан матриц}
\label{tab:competitors-feature}
\end{table}

\textbf{Дүн шинжилгээ.} Хүснэгтээс харахад: (1) Барууны системүүд (Mint, YNAB, Monarch) нь Open Banking-д хүчтэй боловч \emph{Монгол хэл болон Монголын банкны хуулгын форматыг дэмждэггүй}. (2) Wallet, Money Lover нь олон валюттай ажилладаг ч AI-д суурилсан зөвлөгөө байхгүй. (3) Хамгийн алдартай Mint нь 2024 онд хаагдсан тул шинээр хайгчдад орон зай үүссэн. (4) FinTrack-ийн өвөрмөц байр суурь нь --- \textbf{LLM AI чат + Монголын банкны хуулгын автомат AI шинжилгээ + Монгол хэлний бүрэн дэмжлэг} гэсэн гурван онцлогийн нэгдэл бөгөөд эдгээрийг хамтад нь санал болгож буй өөр систем дэлхий дээр одоогоор олдоогүй.

\subsection{Академик судалгааны тойм}
LLM-д суурилсан хувийн санхүүгийн зөвлөгөөний сэдвээр сүүлийн жилүүдэд хийгдсэн судалгааны үндсэн чиглэлүүд:

\begin{itemize}
    \item \textbf{Хэрэглэгчийн итгэлийн судалгаа.} \cite{chatbotfin} өгүүлэлд AI санхүүгийн чатботод хэрэглэгчдийн итгэх хандлагыг шинжилж, тайлбарлан үндэслэх (explainable) хариулт өгөх нь итгэлийг 38\%-иар нэмэгдүүлдгийг тогтоосон.
    \item \textbf{Зан төлвийн ангилал.} PwC \cite{pwc_ai_finance}-ийн судалгаагаар санхүүгийн байгууллагуудын 73\% нь AI-г нэвтрүүлж байгаа боловч \emph{хэрэглэгчид зориулсан} (B2C) шийдлүүд хязгаарлагдмал хэвээр.
    \item \textbf{LLM-ийн контекстэд тулгуурласан хэрэглээ.} \cite{instructgpt, gpt3}-т RLHF болон контекст дотор суралцах хандлагын хүчийг харуулсан нь хэрэглэгчийн санхүүгийн өгөгдлийг шууд зааварчилгаанд оруулж зөвлөгөө гаргах боломжийг нээсэн.
    \item \textbf{Хуулга задлах OCR / NLP.} Уламжлалт OCR + правил-д суурилсан задлах систем нь форматын өөрчлөлтөд эмзэг байсныг \cite{aifinance}-т анхааруулсан бөгөөд LLM-д тулгуурласан хандлага нь илүү уян хатан гэдгийг батлав.
\end{itemize}

Эдгээр судалгаа нь бүгд \emph{Англи хэлний орчинд} хийгдсэн бөгөөд Монгол хэл, Монголын банкны хуулгын форматанд тохирсон тусгайлсан судалгаа байхгүй байна. Энэ нь энэхүү дипломын ажлын академик орон зайг тодорхойлж байна.

Харин Монгол Улсын зах зээлд хувийн санхүүгийн мэдээллийг олон эх үүсвэрээс нэгтгэн боловсруулдаг, AI-д суурилсан цогц шийдэл харьцангуй ховор байна. Зарим банкны гар утасны программ нь энгийн график, зарлагын статистик зэрэг хязгаарлагдмал функцтэй боловч хэрэглэгчийн зан төлөвт суурилсан дүн шинжилгээ хийх, ухаалаг зөвлөмж өгөх түвшинд хүрээгүй байна.\\

Иймээс энэхүү дипломын ажлын хүрээнд Монгол Улсын нөхцөлд тохирсон, хиймэл оюун ухаанд суурилсан хувийн санхүүгийн гар утасны системийн эх загварыг боловсруулах зорилго тавьсан болно.

\section{Хиймэл оюун ухаан ба санхүүгийн салбар}
\textbf{Хиймэл оюун ухаан} (Artificial Intelligence, цаашид AI гэх) нь компьютерийн системийг хүний оюуны үйл ажиллагааг дуурайн гүйцэтгэх чадвартай болгох зорилготой шинжлэх ухаан, технологийн салбар юм. Үүнд суралцах, мэдээлэл боловсруулах, дүгнэлт хийх, шийдвэр гаргах, хэл яриаг ойлгох зэрэг үйлдлүүд багтдаг.\\

AI технологи нь сүүлийн жилүүдэд санхүүгийн салбарт өргөн хэрэглэгдэх болсон бөгөөд өгөгдлийн шинжилгээ, эрсдэлийн үнэлгээ, хэрэглэгчийн зан төлөвийн судалгаа, автомат зөвлөмжийн систем зэрэг олон чиглэлд ашиглагдаж байна.\\

Хиймэл оюун ухааны хөгжлийн түүхийг ерөнхий байдлаар дараах дөрвөн үе шатанд хуваан авч үздэг:

\begin{table}[ht]
\centering
\footnotesize
\renewcommand{\arraystretch}{1.2}
\rowcolors{2}{softgrey}{white}
\begin{tabular}{|c|p{3cm}|p{7.5cm}|}
\hline
\rowcolor{darknavy}
\textcolor{white}{\textbf{Үе шат}} & \textcolor{white}{\textbf{Огноо}} & \textcolor{white}{\textbf{Гол үйл явдлууд}} \\ \hline
1 & 1950--1970   & Symbolic AI, Тюрингийн тест, Lisp хэл, perceptron \\ \hline
2 & 1980--1990   & Expert system, fuzzy logic, neural network-ийн анхны үе \\ \hline
3 & 2000--2015   & Statistical ML, SVM, Random Forest, Big Data, deep learning эхлэл \\ \hline
4 & 2017-оноос   & Transformer, GPT, BERT, LLM, олон горимт загварууд \\ \hline
\end{tabular}
\caption{AI-ийн хөгжлийн үндсэн үе шатууд}
\label{tab:ai-history}
\end{table}

\subsection{Машин сургалтын үндсэн төрөл}
\textbf{Машин сургалт} нь AI-ийн тулгуур салбар бөгөөд өгөгдлөөс автоматаар хууль олох тооцооллын аргуудыг агуулдаг. Гурван үндсэн төрөлд хуваагдана:
\begin{itemize}
    \item [1.] \textbf{Хяналттай сургалт.} Шошготой өгөгдөл дээр сургана. Жишээ: гүйлгээний ангилал (label = "хүнс").
    \item [2.] \textbf{Хяналтгүй сургалт.} Шошгогүй өгөгдөл дээр кластер илрүүлэх. Жишээ: ижил төстэй зан төлөвтэй хэрэглэгчдийг бүлэглэх.
    \item [3.] \textbf{Бататгалын сургалт.} Орчинтой харилцан үйлчлэх замаар хариуцлагын дохиогоор сурах. Жишээ: алгоритмын арилжаа.
\end{itemize}

\subsection{AI-ийн санхүүгийн салбарт хэрэглэгдэх чиглэлүүд}
Санхүүгийн салбар нь AI-ийн хамгийн хүчтэй өсөлттэй салбаруудын нэг юм. PwC-ийн 2024 оны судалгаагаар санхүүгийн байгууллагуудын \textbf{73\%} нь AI-г өөрсдийн үйл ажиллагаандаа нэвтрүүлж байна \cite{pwc_ai_finance, aifinance}. AI-д суурилсан санхүүгийн чатботын хэрэглэгчийн итгэлийн талаарх дэлгэрэнгүй шинжилгээг \cite{chatbotfin} өгүүлэлд танилцуулсан болно.

\begin{table}[ht]
\centering
\footnotesize
\renewcommand{\arraystretch}{1.3}
\rowcolors{2}{softgrey}{white}
\begin{tabular}{|p{3.2cm}|p{4.2cm}|p{4.5cm}|}
\hline
\rowcolor{darknavy}
\textcolor{white}{\textbf{Чиглэл}} & \textcolor{white}{\textbf{Гол асуудал}} & \textcolor{white}{\textbf{AI-ийн оруулах хувь нэмэр}} \\ \hline
Залилан илрүүлэх           & Гүйлгээ хэвийн эсэх       & Anomaly detection, gradient boosting \\ \hline
Зээлийн эрсдэлийн үнэлгээ  & Зээлдэгчийн төлбөрийн чадвар & Logistic regression, deep models \\ \hline
Алгоритмын арилжаа         & Зах зээлийн чиг хандлага  & Time-series загвар, RL \\ \hline
Хэрэглэгчийн шинжилгээ        & Зан төлвийн загвар        & Clustering, classification \\ \hline
Чат хэлбэрийн зөвлөгөө     & Хүн--AI харилцан үйлчлэл  & LLM \\ \hline
Гүйлгээний автомат ангилал & Текст ангилал         & NLP, LLM tagging \\ \hline
Анхааруулга, дүгнэлт       & Шинэ дүгнэлт гаргах       & LLM summarisation, RAG \\ \hline
Робо-зөвлөх (robo-advisor) & Хөрөнгө оруулалтын зөвлөгөө & Portfolio optimisation \\ \hline
\end{tabular}
\caption{AI-ийн санхүүгийн салбарт хэрэглэгдэх үндсэн чиглэлүүд}
\label{tab:ai-finance}
\end{table}

\subsection{Уламжлалт ML-ийн санхүүгийн зөвлөмжид буй хязгаарлалт}
Хувь хүний санхүүгийн зөвлөмжийн системийг бүтээхэд classical ML аргачлал ашиглах боломжтой ч дараах хүндрэлтэй тулгардаг:
\begin{itemize}
    \item Их хэмжээний шошготой сургалтын өгөгдөл шаарддаг (хэрэглэгч бүрд тусдаа загвар сургах нь хэцүү),
    \item Шинж чанарыг (feature engineering) гараар бүтээх шаардлагатай,
    \item Хэрэглэгчийн чөлөөт асуултад хариулах боломжгүй,
    \item Контекстээ ойлгож тайлбар гаргах чадвар сул,
    \item Шинэ ангилал, шинэ хэрэглээний загвар нэмэгдэх тутамд дахин сургах хэрэгтэй.
\end{itemize}

\subsection{LLM-д суурилсан хандлагын давуу талууд}
Орчин үеийн \textbf{Том хэлний загвар} (Large Language Model, LLM) нь дараах давуу талуудтай:
\begin{itemize}
    \item Хэрэглэгчийн санхүүгийн контекстийг (үлдэгдэл, гүйлгээ, ангилал) шууд зааварчилгаанд нэмж зөвлөгөө авдаг,
    \item Чөлөөт хэлбэрийн асуулт-хариулт боломжтой,
    \item Олон төрлийн данс, валют, нөхцөл байдалд адилхан ажилладаг,
    \item Шинэ гүйлгээ нэмэгдэхэд дахин сургах шаардлагагүй,
    \item Хариултыг хүний хэлээр тайлбарласан байдлаар буцаадаг.
\end{itemize}

\section{Том хэлний загвар (LLM) ба Google Gemini}
\textbf{Том хэлний загвар} нь маш том хэмжээний текст өгөгдөл дээр сургагдсан, контекстээ ойлгож, орчуулан, заавар дагах (instruction-following) хийх чадвартай гүн сургалтын загвар юм \cite{gpt3, instructgpt}. Хүний цуглуулсан мэдлэгийг нэгтгэж параметр болгон хадгалах энэхүү онол нь Devlin нарын BERT \cite{bert}-ээс эхлэлтэй.\\

Хиймэл оюун ухааны үндсэн ойлголтуудтай монгол хэлээр танилцахад \cite{mn_aimongol} ном тохиромжтой хэрэгсэл болдог.

\subsection{Transformer архитектур}
2017 онд Vaswani нар "Attention Is All You Need" \cite{transformer} өгүүлэлд танилцуулсан \textbf{Transformer} архитектур нь нийт орчин үеийн LLM-уудын суурь юм. Гол санаа нь \textbf{self-attention} механизм бөгөөд оролтын дараалал дахь үг бүр бусад үгстэй хэрхэн "хамаатай" гэдгийг автоматаар сурдаг.\\

Transformer-ийн гол давуу талууд:
\begin{itemize}
    \item RNN шиг дараалсан тооцоо хийдэггүй тул GPU дээр зэрэгцээ ажиллана,
    \item Урт контекст ойлгох чадвартай (зөвхөн ойролцоо үгсээ биш бүх үгсийг харьцуулна),
    \item Параметрийн хэмжээг нэмэхэд гүйцэтгэл нь тогтмол сайжирна (\emph{scaling law}),
    \item Pretrain → fine-tune → instruction-tune гэсэн 3 шатлалын аргачлалд тохиромжтой.
\end{itemize}

\subsection{LLM-ийн хувьслын үе}

\begin{table}[ht]
\centering
\footnotesize
\renewcommand{\arraystretch}{1.25}
\rowcolors{2}{softgrey}{white}
\begin{tabular}{|c|p{2.7cm}|c|c|p{4.3cm}|}
\hline
\rowcolor{darknavy}
\textcolor{white}{\textbf{Он}} & \textcolor{white}{\textbf{Загвар}} & \textcolor{white}{\textbf{Параметр}} & \textcolor{white}{\textbf{Контекст}} & \textcolor{white}{\textbf{Гол шинэчлэл}} \\ \hline
2018 & GPT-1            & 117M  & 1024 & Pretrain + fine-tune \\ \hline
2019 & GPT-2            & 1.5B  & 1024 & Zero-shot generalisation \\ \hline
2020 & GPT-3            & 175B  & 2048 & Контекст дотор суралцах \\ \hline
2022 & ChatGPT (GPT-3.5)& 175B  & 4096 & RLHF, чат интерфейс \\ \hline
2023 & GPT-4 / Llama 2  & 1.7T  & 32K  & Multi-step reasoning \\ \hline
2024 & Gemini 1.5 / GPT-4o & --  & 1M   & Multimodal, том контекст \\ \hline
2025 & Gemini 2.0 Flash / Claude 3.5 & -- & 1M+ & Тоол-чадварлаг агент \\ \hline
\end{tabular}
\caption{Хүчирхэг LLM-уудын хувьсал}
\label{tab:llm-evolve}
\end{table}

\subsection{LLM сонгох харьцуулалт}

\begin{table}[ht]
\centering
\footnotesize
\renewcommand{\arraystretch}{1.3}
\rowcolors{2}{softgrey}{white}
\begin{tabular}{|p{2.6cm}|p{2.6cm}|c|c|c|}
\hline
\rowcolor{darknavy}
\textcolor{white}{\textbf{Загвар}} & \textcolor{white}{\textbf{Үйлдвэрлэгч}} & \textcolor{white}{\textbf{Контекст}} & \textcolor{white}{\textbf{Free tier}} & \textcolor{white}{\textbf{Multimodal}} \\ \hline
GPT-4o                  & OpenAI    & 128K       & $\times$      & $\checkmark$ \\ \hline
Claude 3.5 Sonnet       & Anthropic & 200K       & $\times$      & $\checkmark$ \\ \hline
\textbf{Gemini 2.0 Flash} & \textbf{Google} & \textbf{1M} & $\checkmark$ & $\checkmark$ \\ \hline
Llama 3.1 70B           & Meta      & 128K       & өөрөө байршуулалт    & $\times$ \\ \hline
DeepSeek R1             & DeepSeek  & 128K       & $\checkmark$  & $\times$ \\ \hline
\end{tabular}
\caption{2025 оны LLM сонгох харьцуулалт}
\label{tab:llm-cmp}
\end{table}

\subsection{Google Gemini API ба олон нийлүүлэгчийн архитектур}
Google-ийн \textbf{Gemini} нь зураг, видео, аудио, текстийг нэгэн зэрэг ойлгодог олон горимт загвар юм. Дипломын ажилд \texttt{gemini-2.5-flash} загварыг үндсэн сонголт болгон ашигласан. Сонголтын гол шалтгаан:
\begin{itemize}
    \item Үнэгүй түвшинтэй (диплом, прототип шатанд тохиромжтой),
    \item Контекст 1М токен хүртэл (банкны хуулгыг бүхэлд нь оруулах боломжтой),
    \item Хариу буцах хурд өндөр (Flash хувилбарт 1--2 секундийн дотор хариулна),
    \item Go хэлэнд албан ёсны SDK байгаа: \code{github.com/google/generative-ai-go},
    \item Зураг, PDF файлыг шууд оруулах боломжтой (олон горимт).
\end{itemize}

\textbf{Газарзүйн хязгаарлалтын асуудал.} Google Gemini API нь Монгол Улсаас шууд хандах боломжгүй тул үйлдвэрлэлийн орчинд найдваргүй болдог. Үүнийг шийдэхийн тулд систем \textbf{олон нийлүүлэгчтэй} архитектурыг хэрэгжүүлсэн:
\begin{itemize}
    \item \textbf{Үндсэн нийлүүлэгч:} Google Gemini \cite{gemini} (хэрэв ажиллах боломжтой бүс нутгаас хандвал).
    \item \textbf{Нөөц нийлүүлэгч:} \textbf{OpenRouter} \cite{openrouter} --- нэг API түлхүүрээр олон LLM-д OpenAI-ийн нийцтэй интерфейсээр хандах боломжтой гарц. Үнэгүй түвшинд Llama 3.3 70B \cite{llama3}, Qwen 2.5 72B, Gemma 2, Mistral 7B, DeepSeek Chat зэрэг загвар ашиглах боломжтой.
    \item \textbf{Автомат нөөц шилжилт.} Үндсэн нийлүүлэгч 429/5xx алдаа буцаавал, эсвэл geo-block илрүүлбэл (HTTP 451 буюу "User location is not supported" алдаа), систем автоматаар OpenRouter руу шилжинэ.
    \item \textbf{Үнэгүй загварын жагсаалтыг автоматаар татах.} OpenRouter-ын үнэгүй загварын нэр тогтмол солигддог тул систем 1 цаг тутамд \texttt{/api/v1/models} API цэгээс одоогийн идэвхтэй жагсаалтыг автоматаар татна.
\end{itemize}

Энэ архитектур нь Монгол хэрэглэгчдэд тусгайлан анхаарч хийгдсэн бөгөөд бусад үнэгүй LLM нийлүүлэгчийг (Groq, Together AI гэх мэт) ирээдүйд адил хандлагаар нэмэх боломжтойгоор бичигдсэн.

\subsection{Зааварчилгаа боловсруулалт ба контекстэд тулгуурласан зөвлөмж}
\textbf{Зааварчилгаа боловсруулалт} гэдэг нь том хэлний загвараас чанартай хариулт авахын тулд оролтын зааварчилгааг зорилго чиглэлтэйгээр, бүтэцтэй бичих ур чадвар юм. Хувийн санхүүгийн зөвлөмжийн систем нь \textbf{контекстэд тулгуурласан зааварчилгаа} зарчмаар ажиллана. Хэрэглэгчийн асуулт ирэх бүрд:
\begin{enumerate}
    \item Өгөгдлийн сангаас тухайн хэрэглэгчийн \textbf{сүүлийн 30 хоногийн гүйлгээ}, нийт үлдэгдэл, ангиллын задаргааг цуглуулна,
    \item Энэ контекстийг JSON хэлбэрээр системийн зааварчилгаанд оруулна,
    \item Хэрэглэгчийн чөлөөт асуултыг хэрэглэгчийн зааварчилгаа болгож нэмнэ,
    \item LLM API руу илгээж бүтэцтэй хариулт буцаана,
    \item Хариултыг өгөгдлийн санд хадгалж UI-д харуулна.
\end{enumerate}

\begin{figure}[ht]
\centering
\begin{tikzpicture}[
    node distance=1.4cm,
    box/.style={rectangle, rounded corners, draw=darknavy, fill=accentblue!15, minimum width=2.4cm, minimum height=1.1cm, align=center, thick, font=\footnotesize},
    boxdb/.style={box, fill=accentorange!25},
    boxctx/.style={box, fill=accentteal!25},
    boxai/.style={box, fill=accentpurple!25},
    boxout/.style={box, fill=accentgreen!25},
    arrow/.style={-{Stealth[length=3mm]}, very thick, draw=darknavy}
]

\node[box]    (q)   {Хэрэглэгчийн\\асуулт};
\node[boxdb,  below=of q]   (db)  {Гүйлгээ,\\үлдэгдэл\\(өгөгдлийн сан)};
\node[boxctx, right=2.2cm of q] (ctx) {Контекст\\нийлүүлэгч};
\node[boxai,  right=2.2cm of ctx] (ai) {LLM\\хариулагч};
\node[boxout, below=of ai]  (ans) {Бэлэн\\зөвлөмж};

\draw[arrow] (q.east)  -- (ctx.west);
\draw[arrow] (db.east) -| (ctx.south);
\draw[arrow] (ctx.east) -- (ai.west);
\draw[arrow] (ai.south) -- (ans.north);

\end{tikzpicture}
\caption{Контекст-төвт AI зөвлөмжийн ажиллах урсгал}
\label{fig:llm-flow}
\end{figure}

\section{Сервер талын технологи}

\subsection{Go хэл}
\textbf{Go} (Golang) нь Google-ийн 2009 онд танилцуулсан, статик төрөлтэй, хөрвүүлэгддэг, санах ойн хог цэвэрлэгчтэй (garbage-collected) системийн программчлалын хэл юм \cite{golang, donovan}. Robert Griesemer, Rob Pike, Ken Thompson нар бүтээгчид бөгөөд анх C++-ийн нийлмэл байдлаас үүдэлтэй хүндрэлүүдийг багасгаж, илүү энгийн хэлийг бий болгох зорилгоор хөгжүүлжээ.\\

Дараах шалтгааны улмаас энэ хэлийг сервер талд сонгосон:
\begin{itemize}
    \item \textbf{Гүйцэтгэл өндөр.} Машины код руу хөрвүүлэгддэг тул Node.js, Python-оос мэдэгдэхүйц хурдан.
    \item \textbf{Зэрэгцээ ажиллах чадвар (concurrency).} Хэлэнд шингэсэн goroutine болон channel ашиглан хөнгөн урсгал (thread) болон мессеж дамжуулалтыг хялбархан хэрэгжүүлдэг. Goroutine нь 2KB стекээс эхэлдэг тул нэг сервер дээр сая орчим урсгалыг зэрэг ажиллуулах боломжтой.
    \item \textbf{Энгийн дүрэм бичлэг.} Класс, удамшил (inheritance) байхгүй, бүтэц (struct) ба интерфейс ашиглан найдвартай дизайн хийдэг.
    \item \textbf{Статик хоёртын файл.} Эцсийн файл нь нэг бүрэн binary болж гарах тул Docker-т 20МБ-аас бага дүрс үүсгэдэг.
    \item \textbf{Стандарт сан хүчирхэг.} HTTP сервер, JSON, шифрлэлт, өгөгдлийн сангийн \texttt{database/sql} зэрэг үндсэн хэрэгслүүд стандарт сангаас бэлэн ирдэг.
    \item \textbf{Орчин үеийн хэрэгсэл.} \texttt{go fmt} автомат форматлагч, \texttt{go test} нэгдсэн тестийн систем, \texttt{go mod} хамаарлын удирдагч.
\end{itemize}

\subsection{Gin веб хүрээ}
\textbf{Gin} \cite{gin} нь Go-ын HTTP веб хүрээ бөгөөд \texttt{httprouter} дээр суурилсан radix-tree чиглүүлэгчтэй. \texttt{net/http} стандарт сан дээр энгийн API өгдөг боловч өндөр хурдтай ажилладаг. Gin-ийн гүйцэтгэл нь Express.js, Flask-аас 40 дахин хурдан гэж хэмжилтээр батлагджээ. Жишээ:

\begin{lstlisting}[language=Go, basicstyle=\ttfamily\footnotesize, caption={Gin REST API цэг}]
r := gin.Default()
r.POST("/api/v1/transactions", txHandler.Create)
r.GET("/api/v1/dashboard", dashHandler.GetDashboard)
r.Run(":8080")
\end{lstlisting}

Gin-ийн гол боломжууд:
\begin{itemize}
    \item Дундын давхаргын систем (Logger, Recovery, Auth, CORS гэх мэт),
    \item JSON холболт ба баталгаажуулалт,
    \item Тогтмол хариултын формат,
    \item Бүлэг чиглүүлэлт (\texttt{api := r.Group("/api/v1")}),
    \item Параметрт чиглүүлэлт (\texttt{:id}, \texttt{*action}).
\end{itemize}

\subsection{GORM ба PostgreSQL}
\textbf{GORM} \cite{gorm} нь Go-д өргөн ашиглагддаг ORM сан бөгөөд struct-ийг автомат байдлаар хүснэгтэд хөрвүүлж, миграц, асуулга хийх боломжтой. Энэ нь хөгжүүлэгчид ачаалал багатай "code-first" хандлагыг өгдөг.\\

\textbf{PostgreSQL} \cite{postgres, redbook} нь нээлттэй эх үүсвэртэй RDBMS бөгөөд 1996 оноос хойш тогтмол хөгжсөн. ACID шинж, JSON багана, транзакц, foreign key, индекс, GIN, GIST индекс, full-text search зэрэг байгууллагын түвшний боломжуудтай. Өгөгдлийн сангийн системийн зохиомжийн талаар Монгол хэл дээрх тойм мэдээллийг \cite{mn_database} номноос унших боломжтой. Дэлхийн хамгийн том санхүү технологийн компаниуд (Stripe, Robinhood, Square) PostgreSQL-ийг ашигладаг.

\begin{table}[ht]
\centering
\footnotesize
\renewcommand{\arraystretch}{1.25}
\rowcolors{2}{softgrey}{white}
\begin{tabular}{|l|c|c|c|c|}
\hline
\rowcolor{darknavy}
\textcolor{white}{\textbf{Технологи}} & \textcolor{white}{\textbf{Хурд}} & \textcolor{white}{\textbf{Зэрэгцээ ажиллагаа}} & \textcolor{white}{\textbf{Сурахад}} & \textcolor{white}{\textbf{Сонгосон}} \\ \hline
Node.js (Express)   & Дунд      & Event loop  & Хялбар   & $\times$ \\ \hline
Python (FastAPI)    & Удаан     & Async       & Хялбар   & $\times$ \\ \hline
Java (Spring Boot)  & Хурдан    & Thread pool & Хүнд     & $\times$ \\ \hline
\textbf{Go (Gin)}   & \textbf{Маш хурдан} & \textbf{Goroutine} & \textbf{Дунд} & $\checkmark$ \\ \hline
Rust (Actix)        & Маш хурдан & Tokio      & Хүнд     & $\times$ \\ \hline
\end{tabular}
\caption{Сервер талын технологийн харьцуулалт}
\label{tab:be-cmp}
\end{table}

\subsection{JWT баталгаажуулалт}
\textbf{JSON Web Token} (JWT) нь төлөв хадгалахгүйгээр сесс удирдах олон улсын стандарт юм \cite{jwt}. Хэрэглэгч нэвтрэх үед сервер JWT үүсгэн (header.payload.signature) гарын үсэг зурж буцаадаг. Дараагийн хүсэлтэд хэрэглэгч энэ токенийг \texttt{Authorization: Bearer <token>} толгой хэсгээр илгээдэг ба сервер гарын үсгийг дахин шалгана.\\

JWT-ийн давуу талууд:
\begin{itemize}
    \item Төлөвгүй --- сервер дээр сесс хадгалах шаардлагагүй,
    \item Тархмал систем (микросервис) ашиглах нөхцөлд тохиромжтой,
    \item Олон төхөөрөмжөөс зэрэг хандах боломжтой.
\end{itemize}

\subsection{Bcrypt нууц үгийн хэшлэл}
\textbf{Bcrypt} \cite{bcrypt} нь Niels Provos болон David Mazi\`{e}res нарын 1999 онд бүтээсэн нууц үгийн хэш алгоритм юм. Аюулгүй кодын ерөнхий зарчмуудтай OWASP Top 10 \cite{owasp} танилцах нь чухал. Хүчээр таах халдлагад тэсвэртэй болгохын тулд "cost factor" (давталтын хэмжээ) тохируулдаг. Cost = 10 гэдэг нь $2^{10} = 1024$ давтан тооцоолох гэсэн үг.\\

Bcrypt-ийн давуу тал:
\begin{itemize}
    \item Salt нь автомат байдлаар үүсгэгддэг,
    \item Rainbow table халдлагад тэсвэртэй,
    \item GPU-аар ажиллуулахад удаан (дасан зохицох cost).
\end{itemize}

\section{Хэрэглэгч талын технологи}

\subsection{Гар утасны программ хөгжүүлэлтийн өнөөгийн нөхцөл}
Орчин үед гар утасны программ хөгжүүлэхэд 4 үндсэн арга бий:
\begin{enumerate}
    \item \textbf{Native iOS} (Swift / SwiftUI) --- iPhone-д зориулсан, хамгийн хурдан, гэхдээ 2 кодын баазтай.
    \item \textbf{Native Android} (Kotlin / Compose) --- Android-д зориулсан, ижил байдлаар хязгаарлагдмал.
    \item \textbf{Олон платформын native} (React Native, Flutter) --- нэг кодын баазаас iOS, Android руу ажилладаг.
    \item \textbf{Hybrid web} (Ionic, Cordova) --- WebView-д суурилсан, гүйцэтгэл сул.
\end{enumerate}

\subsection{React Native ба Expo}
\textbf{React Native} \cite{reactnative} нь Facebook-ийн 2015 онд нээлттэй болгосон, JavaScript/TypeScript ашиглан iOS, Android-д зориулсан native гар утасны программ хөгжүүлэх хүрээ юм. Гар утасны программ хөгжүүлэлтийн ерөнхий аргачлалын талаар \cite{mn_mobiledev} ном тойм мэдээлэл өгдөг. React Native нь UI-г bridge ашиглан native компонент болгон харуулдаг, өөрөөр хэлбэл WebView ашигладаггүй.\\

\textbf{Expo} \cite{expo} нь React Native-ийн дээд хэсэгт суусан хэрэгслийн багц бөгөөд камер, мэдэгдэл, файл сонгогч, биометр нэвтрэлт зэрэг native API-уудыг хялбараар ашиглах боломж олгоно. Expo-ийн гол давуу талууд:
\begin{itemize}
    \item Xcode, Android Studio суулгахгүйгээр анх удаа эхлүүлэх боломжтой,
    \item OTA (Over-The-Air) шинэчлэлт --- хэрэглэгчдэд App Store-оор дамжуулахгүйгээр шинэчлэлт хүргэх боломж,
    \item EAS Build --- үүлэн дээр iOS, Android бүтээлт хийдэг,
    \item 50+ нэмэлт сан (Expo Camera, Expo Notifications гэх мэт).
\end{itemize}

\subsection{TypeScript}
\textbf{TypeScript} \cite{typescript} нь JavaScript-ийн дээд хэсэгт суусан, статик типтэй хувилбар бөгөөд код хөгжүүлэлтийн үед алдааг компиляцийн үед илрүүлдэг. Microsoft-ийн 2012 онд нээлттэй болгосон энэ хэл нь том хэмжээний JavaScript төсөлд зайлшгүй болсон. TypeScript-ийн давуу талууд:
\begin{itemize}
    \item Refactoring хялбар (нэр өөрчилбөл бүх ашигласан газар автомат шинэчлэгдэнэ),
    \item IDE-ийн autocomplete нь нарийвчлалтай,
    \item API-ийн интерфейс илүү тодорхой (DTO type-аар),
    \item Ажиллах үеийн алдааг бууруулдаг.
\end{itemize}

\subsection{NativeWind v4}
\textbf{NativeWind} \cite{nativewind} нь \textbf{Tailwind CSS} \cite{tailwind}-ийг React Native-д ашиглах боломж олгодог сан юм. Tailwind-ийн utility-first зарчмаар \texttt{className="bg-zinc-900 px-4 py-2 rounded-2xl"} гэж бичсэнээр шууд загварчлал гүйцэтгэнэ. NativeWind v4 нь Babel plugin ашиглан className-уудыг бүтээх үед StyleSheet руу хөрвүүлдэг тул ажиллах үеийн ачаалал багатай.\\

Уламжлалт styled-components-ээс ялгарах гол давуу талууд:
\begin{itemize}
    \item Бичсэн код 60--70\%-аар богино,
    \item Design system-ийг \texttt{tailwind.config.js}-аас төвлөрсөн удирдах,
    \item Харанхуй горимыг \texttt{dark:bg-zinc-950} гэх мэтээр шууд тохируулах,
    \item Hot reload-той бүрэн нийцтэй.
\end{itemize}

\begin{table}[ht]
\centering
\footnotesize
\renewcommand{\arraystretch}{1.25}
\rowcolors{2}{softgrey}{white}
\begin{tabular}{|l|c|c|c|c|}
\hline
\rowcolor{darknavy}
\textcolor{white}{\textbf{Хүрээ}} & \textcolor{white}{\textbf{Хэл}} & \textcolor{white}{\textbf{iOS+Android}} & \textcolor{white}{\textbf{Native UI}} & \textcolor{white}{\textbf{Сонгосон}} \\ \hline
Native iOS (SwiftUI)  & Swift      & $\times$ & $\checkmark$ & $\times$ \\ \hline
Native Android (Compose) & Kotlin  & $\times$ & $\checkmark$ & $\times$ \\ \hline
Flutter               & Dart       & $\checkmark$ & Custom        & $\times$ \\ \hline
\textbf{React Native (Expo)} & \textbf{TS} & $\checkmark$ & $\checkmark$ & $\checkmark$ \\ \hline
Ionic                 & TS         & $\checkmark$ & WebView      & $\times$ \\ \hline
\end{tabular}
\caption{Гар утасны хөгжүүлэлтийн хүрээний харьцуулалт}
\label{tab:fe-cmp}
\end{table}

\section{Банкны хуулга боловсруулах техник}
Хэрэглэгч банкаас татаж авсан хуулга нь дараах форматтай байж болно:
\begin{itemize}
    \item \textbf{PDF формат.} Худалдаа Хөгжлийн Банк, Голомт, Хаан банк зэрэг банкны стандарт формат. Текстийг шууд унших боломжтой эсвэл OCR хэрэгтэй болдог.
    \item \textbf{Excel (xlsx).} Албан ёсны экспорт, баганатай форматтай.
    \item \textbf{CSV.} Энгийн текст, задлан уншихад хамгийн хялбар.
\end{itemize}

\subsection{Хуулга задлан уншигчийн арга техник}
Банкны хуулгыг боловсруулахдаа дараах гурван үндсэн арга техникийг ашиглана:
\begin{enumerate}
    \item \textbf{Шууд текст уншилт.} pdftotext, pdfplumber, PyMuPDF гэх мэт сангуудаар PDF файлын текстийг гаргаж авах.
    \item \textbf{Хүснэгтийн хэв шинжилгээ.} Camelot, Tabula зэрэг сангуудаар PDF файл доторх хүснэгтийн бүтцийг таних.
    \item \textbf{OCR (оптик тэмдэгт таних).} Скан хийсэн PDF болон зургийг Tesseract-аар уншиж текст болгон хувиргах.
\end{enumerate}

\subsection{Хуулга задлан уншигчийн архитектур}
Энэхүү системд хуулга задлан уншигчийг хоёр түвшинд хэрэгжүүлсэн:
\begin{enumerate}
    \item \textbf{Python хэл дээрх задлагч (микросервис).} \texttt{FastAPI} \cite{fastapi} + \texttt{pdfplumber} \cite{pdfplumber} + pandas + openpyxl дээр суурилсан тусдаа үйлчилгээ. Монгол Улсын томоохон банкуудын (Хаан банк, Голомт банк, Худалдаа Хөгжлийн Банк, Хас банк, Төрийн банк) хуулгын форматыг түлхүүр үгэнд тулгуурласан танилтын дүрмээр илрүүлэх чадвартай. Үл мэдэгдэх форматын хувьд ерөнхий мөр сканнер ашиглана.
    \item \textbf{Go хэл дээрх нөөц задлагч.} Python үйлчилгээ ажиллахгүй эсвэл хариу өгөхгүй тохиолдолд хамгийн энгийн зохион байгуулалт бүхий (CSV-тэй ойролцоо) хуулгыг Go хэлээр шууд унших нөөц боломж.
\end{enumerate}
Задлагч үйлчилгээтэй \texttt{POST /parse} API цэгээр \texttt{multipart/form-data} хэлбэрээр харьцдаг. Шинжилгээ хийсний дараа хариулт нь \texttt{ParsedStatement} бүтэц (банкны нэр, орлого, зарлага, хугацаа, эхлэх ба эцсийн үлдэгдэл, гүйлгээний жагсаалт, ангиллын задаргаа) хэлбэрээр буцна.

\subsection{Гүйлгээний автомат ангилал}
Гүйлгээний тайлбар (\textit{description}) талбараас ангиллыг тогтоохын тулд хоёр түвшний хандлага ашиглана:
\begin{enumerate}
    \item \textbf{Дүрэмд суурилсан таних (rule-based).} Тодорхой түлхүүр үг (\texttt{"BOLD"}, \texttt{"GOOGOO"}, \texttt{"SHELL"} гэх мэт) → ангилал хөрвүүлэлт.
    \item \textbf{LLM-д суурилсан нэмэлт таних.} Үлдсэн таниагүй гүйлгээг LLM рүү багц байдлаар илгээж зохих ангиллыг асууж тогтоох.
\end{enumerate}

\section{Контейнержуулалт ба байршуулалт}
\textbf{Docker} \cite{docker} нь программыг контейнер хэлбэрээр багцалж, ажиллах орчныг тогтвортой бүрдүүлдэг технологи юм. Docker-ийн гол санаа нь Linux kernel-ийн \texttt{cgroups} болон \texttt{namespace}-ийн боломжийг ашиглан процессыг тусгаарлах юм. Docker-ийн давуу тал:
\begin{itemize}
    \item Орчны давтамжтай байдал --- "минийх дээр ажиллаж байна" асуудал шийдэгддэг,
    \item Цөөн нөөцтэй ажиллана (VM-аас 10 дахин хөнгөн),
    \item Бичил үйлчилгээ (microservice) архитектурын суурь.
\end{itemize}

\textbf{Docker Compose}-ыг ашиглан олон үйлчилгээг (PostgreSQL + сервер тал + задлагч үйлчилгээ) нэг \texttt{docker-compose.yml} файлаас зэрэг ажиллуулах боломжтой. Энэ нь хөгжүүлэлтийн орчныг хэрэглэгч бүрд ижил байлгах боломжийг хангадаг.

\section{Ашигласан нийт технологийн багц}
Хүснэгт~\ref{tab:stack-summary} нь нийт ашигласан технологийн бүрэн жагсаалтыг харуулна.

\begin{table}[ht]
\centering
\footnotesize
\renewcommand{\arraystretch}{1.25}
\rowcolors{2}{softgrey}{white}
\begin{tabular}{|p{2.7cm}|p{4.3cm}|p{4.6cm}|}
\hline
\rowcolor{darknavy}
\textcolor{white}{\textbf{Давхрага}} & \textcolor{white}{\textbf{Технологи}} & \textcolor{white}{\textbf{Зориулалт}} \\ \hline
Хэрэглэгч тал         & React Native + Expo SDK 55  & Гар утасны UI \\ \hline
Хэрэглэгч талын хэл     & TypeScript                  & Статик төрөлтэй JS \\ \hline
Хэрэглэгч талын загвар   & NativeWind v4 (Tailwind)    & Utility-first загварчлал \\ \hline
Төлөв            & React Context + Hooks       & Auth, харанхуй загвар, өгөгдөл \\ \hline
HTTP клиент      & Axios                       & REST API дуудалт \\ \hline
График сан       & react-native-chart-kit      & Багана, муруй, донат график \\ \hline
Биометрик нэвтрэлт & expo-local-authentication & Face ID / Touch ID нэвтрэлт \\ \hline
Олон хэлний дэмжлэг & i18n-js                  & Монгол / Англи хэл солих \\ \hline
Сервер тал          & Go 1.25 + Gin               & REST сервер \\ \hline
ORM              & GORM v1.25.9                & DB struct хөрвүүлэлт \\ \hline
Баталгаажуулалт  & JWT (HS256) + bcrypt        & Токенд суурилсан + нууц үгийн хэш \\ \hline
Өгөгдлийн сан    & PostgreSQL 16               & Реляц мэдээллийн сан \\ \hline
Обьект хадгалалт & MinIO (S3-нийцтэй)          & Хэрэглэгчийн PDF/XLSX/CSV банкны хуулга, профайл аватар зургийг bucket-д хадгалах \\ \hline
AI үндсэн        & Google Gemini 2.5 Flash      & LLM зөвлөмж, ангилал \\ \hline
AI нөөц          & OpenRouter (Llama 3.3, Qwen 2.5, ...) & Газарзүйн хязгаарлалт үүсэх үед автомат нөөц шилжилт \\ \hline
Хуулга задлагч  & Python 3.11 + FastAPI + pdfplumber & Банкны хуулга задлан унших \\ \hline
PDF унших (Go)   & ledongthuc/pdf              & Задлагч серверт хүрэх боломжгүй үед нөөц PDF унших \\ \hline
Excel унших (Go) & xuri/excelize v2            & XLSX хуулга унших \\ \hline
Цаг хуваарь      & Go goroutine + time         & Өдөр тутмын AI дүгнэлтийн хуваарьчлагч \\ \hline
Контейнер        & Docker + Docker Compose     & Орчны изоляци \\ \hline
Push мэдэгдэл    & Expo Push                   & Өдөр тутмын дүгнэлт \\ \hline
\end{tabular}
\caption{Системд ашигласан нийт технологийн стек}
\label{tab:stack-summary}
\end{table}

\textbf{Дүгнэлт.} Энэ бүлэгт хувийн санхүүгийн автоматжуулалтын хэрэгцээ, дэлхийн зах зээлийн өрсөлдөөн, AI-ийн санхүүгийн салбарт буй хэрэглээ, том хэлний загварын архитектур, LLM API-ийн ерөнхий боломж болон дипломын ажилд сонгосон технологийн багц, тэдгээрийг сонгосон шалтгааныг судлав. Дараагийн бүлэгт энэхүү онолын суурин дээр системийн шаардлага, хэрэглээний тохиолдол, өгөгдлийн загварыг тодорхойлно.
